\documentclass[12pt]{article}
\usepackage[utf8]{inputenc}
\usepackage[T1]{fontenc}
\usepackage{hyperref}
\title{Tamasheq Language Resources for K-12 Curriculum Development}
\author{Compiled by Claude Sonnet 5 (claude-sonnet-5) for Andr\'as Kornai}
\date{}
\begin{document}
\maketitle
\section{Introduction}
Tamasheq (ISO 639-3 \texttt{taq}) is a Berber (Afro-Asiatic) language of Mali and Burkina Faso, part of the Tamashek macrolanguage. It is written in Latin script and in Tifinagh. Only one search was run, so this report is limited.
\section{Primary Educational Resources}
Search summaries say Tamasheq is taught in primary grades and in adult literacy classes in Mali. No named textbooks were found. A study of convergent (Tamasheq--French) pedagogy in Mali was found: \url{https://periodicos.ufcat.edu.br/index.php/lep/article/download/45813/22561/192380}
\section{Linguistic and Reference Materials}
\begin{itemize}
\item Sudlow, \emph{Dictionary of the Tamasheq of North-East Burkina Faso}, which is largely valid for Mali and West Niger: \url{https://www.koeppe.de/titel_dictionary-of-the-tamasheq-of-north-east-burkina-faso}
\item Omniglot overview of the writing systems: \url{https://www.omniglot.com/writing/tamasheq.htm}
\item Wikipedia: \url{https://en.wikipedia.org/wiki/Tamasheq_language}
\end{itemize}
\section{Religious and Community Resources}
None found in these searches.
\section{Digital and Audio Resources}
None found in these searches.
\section{Curriculum Development Context}
Tamasheq is already taught as a subject in some Malian schools, which gives a starting point. The materials used there were not identified, so they would need to be located, for example through the Malian education authorities.
\end{document}